\documentclass[11pt,a4paper]{article}
\usepackage[T1]{fontenc}
\usepackage{lmodern,amsmath,amssymb,graphicx,booktabs,longtable}
\usepackage[margin=22mm]{geometry}
\usepackage[hidelinks]{hyperref}
\usepackage{fancyhdr}
\pagestyle{fancy}\fancyhf{}\lhead{Fixed $m_G L$: $L=aN_t$}\rhead{Canonical campaign logbook}\cfoot{\thepage}
\setlength{\headheight}{14pt}\setlength{\parindent}{0pt}\setlength{\parskip}{6pt}\setlength{\emergencystretch}{2em}
\newcommand{\Ensembles}{41}
\newcommand{\Chains}{168}
\newcommand{\TotalSweeps}{1,648,000}
\newcommand{\RetainedSweeps}{1,556,340}
\newcommand{\UsedSweeps}{1,547,750}
\newcommand{\NumericalMasses}{41}
\newcommand{\TuningPoints}{10}
\newcommand{\ContinuumTargets}{3}
\newcommand{\QHalfAll}{0.16331}
\newcommand{\QHalfDrop}{0.15971}
\newcommand{\QOneAll}{0.26225}
\newcommand{\QOneDrop}{0.25572}
\newcommand{\QTwoAll}{0.45733}
\newcommand{\QTwoDrop}{0.41509}

\newcommand{\plotpage}[3]{\clearpage\section*{#1}\includegraphics[width=\linewidth]{#2}\par #3}
\begin{document}
\begin{center}{\LARGE Fixed $m_G L$ campaign}\\[6pt]{\large Theory, methods and preliminary results}\end{center}

\section*{Current status and scope}
The canonical analysis contains \Ensembles{} physical ensembles and \Chains{}
independent chains: \TotalSweeps{} measured sweeps, \RetainedSweeps{} after
independent thermalization cuts, and \UsedSweeps{} used in complete blocks.
There are \NumericalMasses{} numerical mass estimates and \TuningPoints{}
usable preliminary coupling interpolations. All three targets, $c=0.5,1,2$,
have $N_t=2,3,4$; $c=1$ also has $N_t=6$. Every new bootstrap operation uses
exactly 500 replicas. These counts describe the analyzed register; new
simulation batches are tracked separately below. Raw histories remain on qbig.

The register now also includes three $N_t=1$ matched-coupling ensembles:
twelve independent 8,000-sweep chains completed on lnode19 A100 GPUs on
14 October. Their couplings are the three current $N_t=2$ tuning estimates.
For these rows, the folder's $c_{\rm ref}$ denotes that source target;
the actual $N_t=1$ mass is measured without imposing $m_GN_t=c_{\rm ref}$.
They are excluded from fixed-$c$ interpolation and continuum fits.

Four matched-$y$ ensembles, $y=\beta/N_t$, add sixteen 8,000-sweep chains
completed on 16 October. Three extend $y=1.5$ to $N_t=3,4,6$; the fourth
checks $N_s=160$ against 96 at $N_t=4,\beta=6.2$. These rows also have
\texttt{target\_c=null}, with $c_{\rm ref}=0.5$ identifying their source
family. They enter correlator and scaling comparisons. Fixed-$c$ tuning
and its continuum inputs remain unchanged.

The four 80,000-sweep chains at $c=0.5,N_t=3,N_s=80,\beta=4.4$ have the
explicit flag \texttt{length\_outlier}. Their length differs from the
8,000-sweep standard; the flag does not establish a physics bias. They are
pooled with four independent 8,000-sweep chains of identical parameters.
Eight 10,000-sweep chains are explicitly labelled prior inputs from the
fixed-$\beta/N_t$ campaign, at $N_t=2,\beta=3.0$ and $2.1$.

All defined estimates remain available for visual review. Flags identify
window drift, incomplete error plateaus, limited block counts, chain
differences, numerical covariance problems and interpolation sensitivity.
They propagate to preliminary continuum results. The editable
\path{analysis/mg_nt_fixed/review.csv} records manual window choices and
inclusion decisions. No flag is silently converted into a rejection.

\section*{How to read the results}
Each ensemble lives at
\path{analysis/mg_nt_fixed/c_<c>/Nt_<Nt>/Ns_<Ns>/beta_<beta>/}.
It contains correlators, all fit windows, chain statistics and compact
bootstrap arrays. Its \texttt{plots/} directory shows correlator fits,
residuals and the full mass-window scan; \texttt{plots/first\_stage/}
shows every chain's thermalization cut, bootstrap error versus block size,
the selected block and block counts. Blocking examples are kept outside
this PDF. Comparisons sit at the nearest common physics level.
The CSV column guide is \path{docs/mg_nt_fixed/csv_columns.md}.
Figures retain axis labels and legends, with no displayed titles or footer
notes. Their paths, CSVs and this logbook supply the physics parameters and
diagnostic context. The top-level correlator, ensemble, chain, autocorrelation
and fit-window tables include every analyzed simulation, including both untuned
diagnostic batches.

\clearpage\section*{Theory: what is being held fixed?}
The model is three-dimensional SU(2) Yang--Mills theory with Wilson action
\[
 S_W=-\frac{\beta}{2}\sum_p\operatorname{Re}\operatorname{Tr}U_p
       +\text{constant}.
\]
Both transverse directions are periodic.
Tangential links on two temporal Dirichlet surfaces are fixed to the identity
\cite{nebbia}. The separation is $L=aN_t$; the transverse box size is
$L_s=aN_s$. Current simulations store $N_t+1$ spatial-link slices and
$N_t$ temporal-link layers, with no closing layer. Thus $N_t$ counts links
between the fixed surfaces. Historical wrapping-sector outputs retain
their original executable provenance. With fixed surface links their
action factorizes into the physical slab and an independent auxiliary
sector, so removing that sector preserves physical slab expectation values.

For integer transverse separation $r$, define
\begin{align*}
 G(\mathbf{x})&=\prod_{t=0}^{N_t-1}U_t(\mathbf{x},t),\\
 C_G(r)&=\left\langle\tfrac12\operatorname{Re}\operatorname{Tr}
   [G(\mathbf{x})G^\dagger(\mathbf{x}+r\hat e)]\right\rangle.
\end{align*}
Origins and both transverse axes are averaged. The normalization gives
$C_G(0)=1$. The residual boundary transformations act as
$G\mapsto V_LGV_R^\dagger$; in $GG^\dagger$, the $V_R$ factors cancel and
the trace removes the remaining $V_L$ conjugation. This is the
$\mathrm{SU}(2)_L\times\mathrm{SU}(2)_R$ symmetry of the principal chiral
model (PCM). At $N_t=1$ the slab plaquettes reduce exactly to nearest-neighbor
PCM interactions \cite{nebbia}.

The asymptotic decay of $C_G$ measures the lightest transverse screening
scale with overlap in this holonomy channel, under single-state dominance.
With $\widehat m=am_G$, the trajectory is
\[
 c=m_GL=\widehat m N_t\in\{0.5,1,2\},\qquad
 \xi_{\exp}/a=1/\widehat m=N_t/c,\qquad \xi_{\exp}/L=1/c.
\]
The correlation length is twice, equal to, or half the slab thickness.
These are geometry-to-correlation-length ratios; identifying a phase or a
limiting effective theory requires additional evidence. The smallest target
has only a factor-two length separation. A controlled PCM long-distance
interpretation would benefit from $m_GL\ll1$ and agreement of further
dimensionless observables. The present scope is the holonomy mass.

Increasing $N_t$ at fixed $c$ gives, step by step,
\[
 a/L=1/N_t\to0,\qquad \widehat m=c/N_t\to0,\qquad
 \xi_{\exp}/a=N_t/c\to\infty.
\]
Thus the same physical thickness is increasingly well resolved in lattice
units, using the measured screening length as the ruler. The transverse
volume is separately controlled through $m_GL_s=\widehat mN_s$.

\clearpage\section*{From correlator to continuum mass}
At each $N_t$, nearby couplings are simulated and their masses determine
$\beta_\star$ satisfying $\widehat m(\beta_\star,N_t)N_t=c$.
The tuning condition fixes $m_GL$ by construction. The nontrivial continuum
information is the coupling required to satisfy it. In the isotropic Wilson
convention $\beta=4/(ag_{3,0}^2)$ \cite{teper},
\begin{align*}
 g_{3,0}^2L&=(ag_{3,0}^2)N_t=4N_t/\beta_\star,\\
 Q(N_t)\equiv m_G/g_{3,0}^2
 &=\frac{am_G}{ag_{3,0}^2}
 =\frac{\widehat m\beta_\star}{4}
 =\frac{c\beta_\star}{4N_t}.
\end{align*}
The plots fit $Q(N_t)=Q_0+d/N_t$ only ($p=1$). The equivalent
$\beta_\star/N_t$ plots are in each target's comparison directory.
At finite spacing these are bare-coupling coordinates; interpretation in
continuum coupling units assumes the usual continuum matching. At fixed
boundary prescription and large transverse volume, the result determines
the local inverse relation between $m_GL$ and $g_3^2L$.

The $p=1$ form is a restricted cutoff ansatz. Fits using all resolutions
and omitting the coarsest resolution show sensitivity to this choice.
Two points determine a line with zero residual degrees of freedom and
provide no residual test of the ansatz. Statistical bands are conditional
on the chosen mass windows, interpolation windows and volumes. The CSV also
reports a separate conservative linear-response bound on the continuum
intercept from independently signed two-point interpolation shifts at each
resolution. This bound is not added in quadrature to the statistical error.

\section*{Mass analysis and correlated uncertainty}
Each chain is cut independently with MSER-5 and a 500-sweep minimum;
the largest cut across monitored observables is used. Profiles at cadence
five enter as equilibration monitors. Integrated autocorrelation times use
the self-consistent Sokal window and are reported in sweeps \cite{sokal}.
One composite sweep is a heatbath update and five overrelaxation updates.
Holonomy observables are measured every sweep.
The second-moment length is monitored using ensemble structure factors,
\[
 \xi_{\rm 2nd}/a=\frac{1}{2\sin(\pi/N_s)}
 \sqrt{\frac{\widetilde C_G(0)}{\widetilde C_G(p_{\min})}-1},
 \qquad ap_{\min}=2\pi/N_s.
\]
The same joint replicas propagate its error. Ratios are formed after
ensemble averaging; unresolved replicas are recorded explicitly.

Complete nonoverlapping blocks are formed inside each chain. Candidate
lengths scan the whole post-cut history, starting at five sweeps. The chosen
length exceeds twice the largest measured autocorrelation time and uses
the earliest common three-length error plateau when available. Each curve
point is a separate 500-replica bootstrap SE. Plateau tolerances are 5\%
or two approximate SE uncertainties, whichever is larger; the latter use
Welch--Satterthwaite degrees of freedom and the Monte Carlo SE of 500 draws.
This is a diagnostic stability criterion. Missing plateaus are flagged.
Fewer than 30 complete blocks is a warning; 50 is the target.

\clearpage\section*{Pooling, fitting and window review}
If chain $k$ contributes $n_k$ complete blocks of $B_k$ measurements,
$w_k=n_kB_k/\sum_j n_jB_j$. In each of exactly 500 replicas, draw $n_k$
block vectors with replacement within chain $k$, then form
$C^{(b)}=\sum_k w_k\overline C_k^{(b)}$ \cite{efron}.
The draw acts on the entire separation vector. Chain identities and weights
are fixed. The ensemble mean is $\sum_k w_k\overline C_k$.
Independent chains retain their own cuts, autocorrelations and block
boundaries. Split-history and between-chain discrepancies are reported
separately; maxima over many correlated observables are review diagnostics.

The asymptotic two-image model is \cite[Eq.~(4.14)]{nebbia}
\[
 C_G(r)\simeq A\{K_0(\widehat m r)+K_0[\widehat m(N_s-r)]\}.
\]
Its interpretation requires a stable long-distance window and finite-volume
control. Additional images, excited contributions and lattice dispersion
may matter outside that regime. Each central and bootstrap fit minimizes
the diagonal-error objective. Parameter errors are the SD of the same fit
over the 500 vector replicas. The complete covariance is used to report
$\chi^2_{\rm full}=r^T\operatorname{Cov}^{-1}r$ after fitting.
This statistic is descriptive: the fitted parameters do not minimize it,
so a nominal analytic chi-square p-value is not appropriate. No p-value is
computed and no chi-square gate is imposed. Singular covariance leaves this
diagnostic undefined while retaining a defined diagonal fit.

For the fixed-$c$ and matched-$y$ ensembles, all contiguous windows with at least five points are fitted, with every
positive $R_{\min}$ and $R_{\max}=N_s/2,N_s/2-2,\ldots\geq\lceil N_s/3\rceil$.
The recommendation is the earliest stable $3\times3$ patch in adjacent
endpoints, preferring the largest $R_{\max}$. Pairwise masses must differ
by at most two combined bootstrap SEs. Overlapping windows are correlated;
this rule is a heuristic for visual review. If no such patch exists, the
least-drifting numerical patch is flagged and retained. Failed refit slots
remain NaN in the saved 500-slot arrays; success counts and failure reasons
are in CSV. Relative error above 25\%, low success fraction and a missing
mass-error blocking plateau produce flags.

For the new $N_t=1$ boxes with $N_s=512,2048$, the user approved a bounded
endpoint scan: seven allowed upper endpoints and logarithmically spaced
lower endpoints, followed by dense adjacent checks near three candidate
plateaus. Every integer separation inside each window and all 500 refits
are retained. The recommendation is the earliest stable patch among tested
windows; it does not establish the exhaustive earliest patch. The $N_s=128$
scan remains exhaustive. A covariance formed from 500 centered replicas
has rank at most 499; windows with 500 or more separations therefore have
undefined full-covariance $\chi^2$. Fewer independent block vectors can
reduce the rank further. No inverse regularization is used to manufacture
a goodness-of-fit statistic.

The collector enforces the additional block-count bound
$\operatorname{rank}(\operatorname{Cov})\leq\min[499,\sum_k(n_k-1)]$:
within each chain, centered resampling multiplicities sum to zero, leaving
at most $n_k-1$ independent directions. Wider windows have undefined full
$\chi^2$ even if roundoff allows Cholesky factorization. Diagonal-fit masses
and their 500 refits are unchanged by this diagnostic correction.

\clearpage\section*{Coupling windows and volume choices}
For each $c,N_t$, all contiguous three-coupling windows are examined using
$\log(\widehat mN_t)=a_0+b\beta$. Only centrally bracketed targets yield
an interpolation; undefined slopes or unbracketed targets are recorded.
Each aligned mass replica gives $\beta_\star^{(b)}=(\log c-a_0^{(b)})/b^{(b)}$.
The median and 16--84\% interval are quoted. The largest shift relative to
a target-bracketing two-point interpolation is a separate systematic.
The closest bracketing window is suggested; manual choices remain editable.

The larger $N_s=96$ boxes are primary at $c=1,N_t=4,\beta=4.95,5.3$;
the $N_s=48$ versions remain volume diagnostics. At $c=0.5,N_t=4,
\beta=5.58394149$, $N_s=96$ is primary and 128 is a matched-volume check.
The $N_t=4,c=1$ confirmation ensemble now contributes to tuning. Its earlier
independence applies to the historical interpolation only.

\section*{Coupling and continuum results}
\begin{center}\small\begin{tabular}{rrrrrl}
\toprule
$c$ & $N_t$ & $\beta_\star$ & 16--84\% interval & interp. shift & status\\\midrule
0.5 & 2 & 2.9895 & [2.9888, 2.9903] & 0.00276 & flagged\\
0.5 & 3 & 4.3069 & [4.3059, 4.3081] & 0.000383 & flagged\\
0.5 & 4 & 5.5846 & [5.5821, 5.5868] & 0.000685 & flagged\\
1 & 2 & 2.5045 & [2.5038, 2.5051] & 0.000365 & flagged\\
1 & 3 & 3.5985 & [3.5968, 3.6002] & 0.00201 & flagged\\
1 & 4 & 4.6305 & [4.6279, 4.6326] & 0.00352 & flagged\\
1 & 6 & 6.6664 & [6.6636, 6.6691] & 0.0108 & flagged\\
2 & 2 & 1.9624 & [1.9603, 1.964] & 0.00269 & flagged\\
2 & 3 & 2.9428 & [2.941, 2.9442] & 0.00187 & flagged\\
2 & 4 & 3.773 & [3.7704, 3.7754] & 0.00349 & flagged\\
\bottomrule
\end{tabular}
\end{center}
All entries are preliminary and subject to visual review. The interpolation
shift is kept separate from the statistical interval.
\begin{center}\small\begin{tabular}{rlrrrr}
\toprule
$c$ & resolutions & $Q_0$ & boot. SE & interp. bound & dof\\\midrule
0.5 & 2, 3, 4 & 0.16331 & 0.000124 & 0.000295 & 1\\
0.5 & 3, 4 & 0.15971 & 0.000331 & 0.000134 & 0\\
1 & 2, 3, 4, 6 & 0.26225 & 0.000158 & 0.000647 & 2\\
1 & 3, 4, 6 & 0.25572 & 0.000261 & 0.00108 & 1\\
2 & 2, 3, 4 & 0.45733 & 0.000768 & 0.00163 & 1\\
2 & 3, 4 & 0.41509 & 0.00153 & 0.00268 & 0\\
\bottomrule
\end{tabular}
\end{center}
The second row for each target omits its coarsest resolution. Detailed flags,
full-covariance residual statistics and bootstrap success counts are in the
CSV tables. The chosen primary volumes are not a volume extrapolation.

\textbf{Resolution dependence is the principal limitation of these extrapolations.}
Omitting $N_t=2$ moves $Q_0$ from approximately $\QHalfAll$ to $\QHalfDrop$
at $c=0.5$, $\QOneAll$ to $\QOneDrop$ at $c=1$, and $\QTwoAll$ to
$\QTwoDrop$ at $c=2$. These shifts are much larger than the conditional
bootstrap errors. They are explicitly flagged using the paired-replica
uncertainty of the difference between the two fits. The tables retain all
requested preliminary fits; their small statistical errors do not establish
a controlled continuum mass. Finer resolutions and visual range selection
are the priority before a precision continuum claim.

\section*{Why earlier results changed}
The original resampler included an outer draw of independent chain identities.
The corrected method fixes those identities and resamples blocks within
each chain. Earlier bootstrap counts and uncertainty prescriptions are
preserved in the checksum-verified qbig archive; they are historical products.
This reanalysis additionally pools matching short and long runs with fixed
used-measurement weights and applies the requested 500-replica policy to all
38 physical ensembles.

Earlier acceptance also imposed fit p-value and window gates. In particular,
the $c=2,N_t=3,N_s=32,\beta=3.25$ mass was unavailable after those gates,
despite numerical windows that can be inspected for a plateau. The new
analysis retains numerical windows, flags drift and lets visual review
determine their use. The coupling fits previously rejected for local
log-linearity are displayed in the comparison directories, including their
residuals. Keeping these preliminary results increases coverage; their
statistical errors remain conditional on the model and selected window.

\section*{Next simulations to consider after review}
First inspect flagged mass and coupling windows, especially the $\beta=3.25$
case and the smallest remaining block counts. Additional independent
8,000-sweep chains can improve precision and between-chain checks without
adding another length outlier. At otherwise fixed conditions, doubling
effective statistics reduces statistical error by approximately $\sqrt2$;
halving it requires roughly four times the effective statistics. More sweeps
cannot resolve a poor fit model or uncontrolled finite-volume effects.
At a fixed block size $B$, obtaining 50 blocks within every chain requires
at least its thermalization cut plus $50B$ sweeps: for a 500-sweep cut this
is 10,500 sweeps at $B=200$ and 12,500 at $B=240$. Adding more 8,000-sweep
chains improves combined statistics but leaves that per-chain count unchanged
unless the selected block size changes.

The largest current dataset supports three $c$ values, with four resolutions
only at $c=1$. The next coverage priority is an $N_t=6$ bracket for $c=0.5$
and $c=2$, followed by independent confirmation of the tuned couplings and
volume/timing checks. No new simulation is authorized by this report.

\clearpage\section*{New $N_t=1$ matched-coupling diagnostics}
\begin{center}\small\begin{tabular}{rrrrrrr}
\toprule
$c_{\rm ref}$ & $\beta$ & $N_s$ & $am_G$ & boot. SE & $\xi_{\rm 2nd}/a$ & min. blocks\\\midrule
0.5 & 2.989543 & 2048 & 0.0067043 & 3.77e-05 & 146.7 & 6\\
1 & 2.504456 & 512 & 0.028794 & 6.3e-05 & 40.59 & 36\\
2 & 1.962387 & 128 & 0.14111 & 0.000219 & 7.348 & 37\\
\bottomrule
\end{tabular}
\end{center}
All twelve chains completed with verified inputs, output counts, boundary
conditions and checkpoints, using KLFT revision
\texttt{f47421d895469b72d4af8a9b8479807f2b09a692} and its native A100 build.
Total simulation time was 3.038 A100 GPU-hours. Holonomies were measured
every sweep, profiles every five sweeps, and Wilson loops were disabled.
The separate P100 validation job 290996 also passed all six CTests, CUDA
memory checking and the benchmark; it launched no production.

The transverse sizes 128, 512 and 2048 were explicit candidate input
choices. The saved planning records label them unvalidated and contain no
quantitative derivation establishing their necessity. At the largest
coupling, the first-stage estimate is $\xi_{\rm 2nd}/a\simeq147\pm34$,
giving $N_s/\xi_{\rm 2nd}\simeq14$ at the central value. This is preliminary
because independent cuts range from 825 to 3625 sweeps, the maximum
estimated autocorrelation time is about 317 sweeps, and $B=635$ leaves
only 6--11 complete blocks per chain. The common blocking plateau is
unresolved and split-history/between-chain differences require review.
At $N_s=512$, $\xi_{\rm 2nd}/a\simeq40.6\pm2.9$ and $B=200$ leaves
36--37 blocks; at $N_s=128$, $\xi_{\rm 2nd}/a\simeq7.35\pm0.23$ and
$B=200$ leaves 37 blocks. These measured lengths explain why matched
couplings cannot be interpreted as the source fixed-$c$ targets at $N_t=1$.

The largest-volume statistics are the immediate concern. At the present
block length, 50 complete blocks require 31,750 post-cut sweeps per chain;
the most conservative observed cut would bring this to 35,375 total sweeps.
This is a planning estimate: autocorrelation and thermalization must be
rechecked after extending statistics. A smaller-volume comparison is needed
before replacing 2048 by a smaller production size. Adding these untuned
points does not extend the fixed-$c$ continuum resolution set.

\subsection*{New $N_t=1$ matched-$y$ batch}
On 20 October, the user approved four additional $N_t=1,N_s=128$ ensembles
at $\beta=1.45,1.50,4.6/3,1.55$, each with four independent 8,000-sweep hot
chains. P100 smoke job 291028 and timing array 291029 passed geometry,
boundary, cadence and output-hash checks. The projected cost is 0.438 P100
GPU-hours. Production array 291033 completed all sixteen chains on lnode13
and lnode17 by 22:00:55 CEST: 128,000 sweeps and 0.1525 P100 GPU-hours.
Output counts, input hashes and physical checkpoints passed verification.
This batch has no mass analysis yet and is excluded from the tables and
figures below; its target remains \texttt{null}.

\clearpage\section*{Matched-$y$ extension near $\beta/N_t=1.5$}
\begin{center}\small\begin{tabular}{rrrrrrr}
\toprule
$N_t$ & $N_s$ & $\beta$ & $am_G$ & boot. SE & $m_G L$ & min. blocks\\\midrule
3 & 96 & 4.5 & 0.13865 & 0.000348 & 0.41594 & 93\\
4 & 160 & 6 & 0.094149 & 0.000193 & 0.3766 & 125\\
4 & 160 & 6.2 & 0.082128 & 0.000186 & 0.32851 & 93\\
6 & 320 & 9 & 0.056306 & 0.000126 & 0.33784 & 31\\
\bottomrule
\end{tabular}
\end{center}
Each ensemble has four independent 8,000-sweep hot-start chains, measured
every sweep (profiles every five). The batch used 0.2745 A100 GPU-hours on
lnode19, preserving one free GPU. Inputs, output counts, checkpoints and
raw hashes passed validation. Array 291008 used the same 500-replica
estimator and exhaustive window scan as the fixed-$c$ analysis.

All four mass scans completed successfully. Their selected windows have
500 successful refits. The $N_t=3$ ensemble has a split-history warning;
$N_t=4,\beta=6$ has between-chain and later-window drift warnings. The
$N_t=6$ window is $r=6$--160, with later-window drift retained for review.
Its $B=240$ leaves 31 blocks per chain. The mass-error plateau includes
$B=160,200,240$, while the common observable plateau remains unresolved.
At the current 500-sweep cut, 50 blocks would require 12,500 total sweeps per
chain, followed by renewed equilibration and autocorrelation checks.
The full-covariance statistic is undefined because its rank is at most
120 for 155 fitted separations. The defined diagonal mass remains included.
At $N_t=4,\beta=6.2$, the $N_s=160$ minus 96 mass difference is
$-0.0003541\pm0.0003835$, or 0.92 bootstrap standard deviations, conditional
on the chosen windows.

Visual review of the new $N_t=6$ correlator shows a positive, flattening tail
beyond roughly $r=90$, above the selected two-image $K_0$ curve. Later
lower-endpoint choices pull the fitted mass downward as its error grows.
The origin of that tail remains unresolved. The quoted mass therefore
requires window and equilibration review alongside increased statistics;
the numerical result and every alternative window are retained.

The purpose is to test the proposed collapse where the screening mass is
small, and to check transverse-volume sensitivity. The existing
$N_t=2,N_s=64,\beta=3$ point joins the three new $y=1.5$ points. From the
Wilson coupling convention,
\[
 y=\beta/N_t=\frac{4}{ag_{3,0}^2N_t}
   =\frac{4}{g_{3,0}^2L}.
\]
Thus fixed $y$ holds the bare slab-coupling coordinate fixed while the
resolution changes. Measuring $am_GN_t$ tests its approach to a finite
physical screening scale. If that limit is nonzero, $am_G$ ultimately
scales as $1/N_t$, up to cutoff and coupling-matching corrections.

The displayed $am_GN_t^{1.5}$ collapse uses the user's trial exponent;
no exponent was fitted. For successive resolutions we also tabulate
\[
 \alpha_{\rm eff}(N_1,N_2)=
   \frac{\log[\widehat m(N_1)/\widehat m(N_2)]}{\log(N_2/N_1)}.
\]
Substitution of $\widehat m=A N_t^{-\alpha}$ gives
$\log[(N_2/N_1)^\alpha]/\log(N_2/N_1)=\alpha$.
For measured finite-spacing masses this is a diagnostic log ratio, with
errors propagated through the saved 500 replicas. Window, volume and
thermalization uncertainties remain separate from these conditional
statistical errors. The fixed-$c$ continuum fit continues to use its
existing independently tuned inputs.
The measured adjacent-resolution powers are $1.4054\pm0.0077$ (2 to 3),
$1.3455\pm0.0111$ (3 to 4), and $1.2678\pm0.0075$ (4 to 6).
Their drift shows that the trial exponent 1.5 is an approximate visual
scaling choice over this range. These finite-spacing ratios inherit the
window and statistics warnings; they do not determine an asymptotic exponent.

\plotpage{Lorenzo's trial collapse with the matched-$y$ batch}{analysis/mg_nt_fixed/comparison/lorenzo_scaling_collapse.pdf}{Reproduction of Lorenzo's supplied $am_GN_t^{1.5}$ versus $\beta/N_t$ plot, including all three $N_t=1$ matched-coupling diagnostics. Filled markers identify the four matched-$y$ ensembles; open markers show earlier mass estimates. The exponent 1.5 is an assumed illustrative choice. All defined flagged masses remain visible, with bootstrap standard errors conditional on their chosen windows.}
\plotpage{Screening mass in slab units at all measured couplings}{analysis/mg_nt_fixed/comparison/mass_Nt_vs_beta_over_Nt.pdf}{The ordinate is $m_GL=am_GN_t$, with $L=aN_t$. Each saved mass and its bootstrap standard error are multiplied by the exact temporal extent. The three analyzed $N_t=1$ ensembles are included at their measured $\beta/N_t$ values, approximately 1.962, 2.504 and 2.990. They retain their untuned diagnostic role. The new $N_t=1,\beta/N_t=1.5$ mass is not yet analyzed, so the matched-coupling comparison on the next page retains its existing resolution set. No new fit or bootstrap resampling is used.}
\plotpage{Measured screening scale at $y=1.5$}{analysis/mg_nt_fixed/c_0p5/comparison/matched_y_mass_Nt.pdf}{The new $N_t=3,4,6$ ensembles extend the earlier $N_t=2$ point. No fixed-$c$ target is imposed on this comparison and no additional continuum fit is performed.}
\plotpage{Second-moment length in slab units}{analysis/mg_nt_fixed/comparison/xi_2nd_slab_scaled.pdf}{All measured couplings are shown, including the $N_t=1$ matched-coupling diagnostics. The logarithmic scale accommodates their larger lengths. The coordinate is $\xi_{\rm 2nd}/L=(\xi_{\rm 2nd}/a)/N_t$; companion figures show its square, unscaled lengths versus $\beta$, and the trial normalization $N_t^{1.5}$. Error bars use the saved 500 joint replicas and retain the volume and statistics limitations. For the squared estimator, evaluate $(\widetilde C(0)/\widetilde C(p_{\min})-1)/[4\sin^2(\pi/N_s)]$ on each replica before taking its SD. Negative noisy squared-moment estimates remain in that distribution. Both central values and errors are in the existing ensemble CSVs.}
\plotpage{Transverse-volume check at $N_t=4,\beta=6.2$}{analysis/mg_nt_fixed/c_0p5/Nt_4/comparison/volume_beta_6p2.pdf}{Independent $N_s=96$ and 160 ensembles are compared using their selected mass windows. The companion CSV reports the difference and its 500-replica uncertainty; other windows and correlator curves remain available for visual review.}

\plotpage{$N_t=1$: matched-coupling correlators}{analysis/mg_nt_fixed/comparison/nt1_correlators.pdf}{The three volumes correspond to different couplings. Their folder labels identify the source $N_t=2$ tuning. Error bars use independent-chain, pooled-with-boundaries bootstrap replicas.}
\plotpage{$N_t=1$: preliminary screening masses}{analysis/mg_nt_fixed/comparison/nt1_masses.pdf}{Conditional bootstrap standard errors are shown. Window, blocking, chain-consistency and volume flags remain in the per-ensemble CSVs. These masses are excluded from fixed-$c$ continuum fits.}

\plotpage{Continuum comparison: all targets}{analysis/mg_nt_fixed/comparison/continuum_all_targets_p1.pdf}{All three targets use the requested $p=1$ ansatz. Flags remain attached to their input points.}
\plotpage{$c=0.5$: resolution sensitivity}{analysis/mg_nt_fixed/c_0p5/comparison/continuum_mass_in_coupling_units_p1.pdf}{The drop-coarsest line uses two resolutions and has zero residual degrees of freedom. Bands contain conditional statistical uncertainty only.}
\plotpage{$c=1$: four measured resolutions}{analysis/mg_nt_fixed/c_1/comparison/continuum_mass_in_coupling_units_p1.pdf}{The available $N_t=6$ bracket extends the continuum comparison. Independent tuning confirmation and further cutoff control remain useful.}
\plotpage{$c=2$: preliminary retained trajectory}{analysis/mg_nt_fixed/c_2/comparison/continuum_mass_in_coupling_units_p1.pdf}{Mass-window and interpolation concerns remain visible in the CSV flags and per-ensemble plots.}
\plotpage{Previously gated mass: $N_t=3,N_s=32,\beta=3.25$}{analysis/mg_nt_fixed/c_2/Nt_3/Ns_32/beta_3p25/plots/mass_windows_Rmax_16.pdf}{Every numerical mass window at this endpoint is shown. The band marks the suggested mass; overlapping-window stability is a diagnostic heuristic. Other endpoints and the correlator residuals are saved alongside this plot.}
\plotpage{Example coupling interpolation}{analysis/mg_nt_fixed/c_2/Nt_3/comparison/tuning.pdf}{The previously gated mass is retained as a preliminary input. All coupling windows and their residual plots are available in the same comparison directory.}

\clearpage\section*{Fixed-$c$ ensemble result register}
\small\begin{longtable}{rrrrrrrl}
\toprule
$c$ & $N_t$ & $N_s$ & $\beta$ & $am_G$ & SE & window & used sweeps\\\midrule\endhead
0.5 & 2 & 64 & 2.9 & 0.28458 & 0.000491 & 3--32 & 29,760\\
0.5 & 2 & 64 & 3 & 0.24513 & 0.000456 & 4--32 & 37,600\\
0.5 & 2 & 64 & 3.1 & 0.21427 & 0.000438 & 3--32 & 29,920\\
0.5 & 3 & 80 & 4.2 & 0.18478 & 0.000372 & 3--40 & 29,440\\
0.5 & 3 & 80 & 4.27 & 0.17315 & 0.000339 & 3--40 & 30,000\\
0.5 & 3 & 80 & 4.29 & 0.16951 & 0.000328 & 3--40 & 30,000\\
0.5 & 3 & 80 & 4.4 & 0.15159 & 0.000156 & 5--40 & 347,000\\
0.5 & 3 & 80 & 4.6 & 0.12704 & 0.000299 & 3--40 & 29,760\\
0.5 & 4 & 96 & 5.5 & 0.1334 & 0.000399 & 4--48 & 29,440\\
0.5 & 4 & 96 & 5.58394149 & 0.125 & 0.000371 & 4--48 & 29,600\\
0.5 & 4 & 96 & 5.6 & 0.12358 & 0.000333 & 4--48 & 29,760\\
0.5 & 4 & 96 & 5.9 & 0.10119 & 0.000325 & 4--48 & 30,000\\
0.5 & 4 & 96 & 6.2 & 0.082482 & 0.000326 & 4--48 & 29,760\\
0.5 & 4 & 128 & 5.58394149 & 0.12562 & 0.000255 & 4--64 & 29,760\\
1 & 2 & 32 & 2.35 & 0.6181 & 0.00139 & 2--16 & 29,920\\
1 & 2 & 32 & 2.5 & 0.50322 & 0.000599 & 1--16 & 30,000\\
1 & 2 & 32 & 2.65 & 0.40906 & 0.000958 & 2--16 & 29,760\\
1 & 3 & 48 & 3.5 & 0.36806 & 0.000895 & 3--24 & 30,000\\
1 & 3 & 48 & 3.7 & 0.3004 & 0.000733 & 3--24 & 29,920\\
1 & 3 & 48 & 3.9 & 0.2481 & 0.00064 & 3--24 & 30,000\\
1 & 4 & 48 & 4.6 & 0.25524 & 0.000661 & 3--24 & 29,920\\
1 & 4 & 48 & 4.61755576 & 0.25314 & 0.000718 & 3--24 & 30,000\\
1 & 4 & 48 & 4.95 & 0.19764 & 0.000578 & 3--24 & 29,760\\
1 & 4 & 48 & 5.3 & 0.15342 & 0.000614 & 3--24 & 29,760\\
1 & 4 & 96 & 4.95 & 0.19631 & 0.00042 & 4--48 & 30,000\\
1 & 4 & 96 & 5.3 & 0.15284 & 0.000385 & 4--48 & 30,000\\
1 & 6 & 96 & 6.4 & 0.19015 & 0.000364 & 4--48 & 29,760\\
1 & 6 & 96 & 6.7 & 0.16289 & 0.000486 & 5--48 & 30,000\\
1 & 6 & 96 & 7 & 0.14187 & 0.000362 & 4--48 & 29,760\\
2 & 2 & 32 & 1.9 & 1.0772 & 0.00403 & 2--16 & 29,440\\
2 & 2 & 32 & 2 & 0.95897 & 0.00322 & 2--16 & 29,920\\
2 & 2 & 48 & 2.1 & 0.83975 & 0.00392 & 3--24 & 37,920\\
2 & 3 & 32 & 2.85 & 0.73663 & 0.00198 & 2--16 & 29,760\\
2 & 3 & 32 & 3.05 & 0.59445 & 0.00143 & 2--16 & 29,920\\
2 & 3 & 32 & 3.25 & 0.47481 & 0.0019 & 3--16 & 30,000\\
2 & 4 & 32 & 3.7 & 0.53286 & 0.00134 & 2--16 & 29,920\\
2 & 4 & 32 & 4 & 0.40814 & 0.0016 & 3--16 & 30,000\\
2 & 4 & 32 & 4.3 & 0.323 & 0.00125 & 3--16 & 29,920\\
\bottomrule
\end{longtable}
\normalsize
The untuned diagnostic registers are in their respective sections above.
Uncertainties are bootstrap standard errors of $am_G$. The selected window
and used statistics are included so a result can be traced to its CSV and
raw-input hashes. The complete chain-level cuts, autocorrelation times,
flags and block counts are in \texttt{chains.csv} and
\texttt{autocorrelations.csv}; no chain histories are concatenated.

\clearpage\begin{thebibliography}{9}
\bibitem{nebbia} M. Nebbia, \emph{Effects of Dirichlet Boundary Conditions in Yang--Mills Theories}, MSc thesis (2024/2025), local reference \path{docs/marconebbia_tesi.pdf}.
\bibitem{teper} M. Teper, SU($N$) Yang--Mills theory in $2+1$ dimensions, Phys. Rev. D 59, 014512 (1999), \href{https://arxiv.org/abs/hep-lat/9804008}{hep-lat/9804008}.
\bibitem{sokal} A. D. Sokal, Monte Carlo Methods in Statistical Mechanics: Foundations and New Algorithms, in \emph{Functional Integration}, NATO ASI B 361 (1997), pp. 131--192.
\bibitem{efron} B. Efron and R. J. Tibshirani, \emph{An Introduction to the Bootstrap}, Chapman \& Hall (1993), chapters 6 and 8.
\end{thebibliography}
\end{document}
